\documentclass[10pt,a4paper]{amsart}
\usepackage[margin=19mm]{geometry}
\usepackage{amsmath,amssymb,mathtools}
\usepackage[scr=rsfs,cal=euler]{mathalfa}
\usepackage{xfrac}
\usepackage[unicode,hidelinks,bookmarks=false]{hyperref}
\newcommand{\ie}{i.e.\ }
\newcommand{\cf}{cf.\ }
\newcommand{\resp}{respectively\ }
\newcommand{\Subsection}{\subsection}
\providecommand{\nobreakdash}{\nobreak\hbox{-}}
\setlength{\emergencystretch}{2em}
\multlinegap0pt
\usepackage{xspace,mathrsfs}
\newtheorem{theorem}{Theorem}[section]
\newtheorem{lemma}[theorem]{Lemma}
\newtheorem{definition}[theorem]{Definition}
\newtheorem{example}[theorem]{Example}
\newtheorem{exercise}[theorem]{Exercise}
\newtheorem{proposition}[theorem]{Proposition}
\newtheorem{corollary}[theorem]{Corollary}
\newtheorem{lem}[theorem]{Lemma}
\newtheorem{note}[theorem]{Note}
\newtheorem{remark}[theorem]{Remark}
\numberwithin{equation}{section}
\newcommand{\abs}[1]{\LVERT#1\RVERT}
\begin{document}
\title[Q-filters of bounded lattices]{Q-filters of bounded lattices}
\author{[2020] 97H50, 06A11 ęywords $ Q $-filter, $S$-$ Q $-filter, almost $S$-$ Q $-filter; Mahdi Anbarloei}
\date{}
\maketitle
\noindent\emph{Source and licence.} Q-filters of bounded lattices. Version arXiv:2606.00731v1 (CC BY 4.0). This material is distributed under the Creative Commons Attribution 4.0 International licence, http://creativecommons.org/licenses/by/4.0/, as recorded in the arXiv OAI licence metadata for this exact version and shown on the abstract page https://arxiv.org/abs/2606.00731v1. Full rights notice: licenses/arxiv-2606.00731-CC-BY-4.0.txt.\par\smallskip
\noindent\textit{Editorial scope.} Counted unit: the original \texttt{theorem} environment \texttt{2} at source lines 175--181 together with its complete proof at lines 182--186; the delivered document keeps source lines 143--195 so that every referenced label and every citation resolves inside it. Native editable source is the paper's own concatenated TeX; the AMS wrapper is editorial formatting only. No publisher class, upstream script or unrelated artwork is redistributed.
\section{$\mathcal{Q}$-filters}
In this section, we introduce and investigate the  concept of $\mathcal{Q}$-filters in a bounded lattice. We initiate our study with the following definition.
\begin{definition}
Let $\mathcal{P}$ and $ \mathcal{Q}$ be  proper filters of $\mathcal{L}$. We say that $\mathcal{P}$ is a {\it $\mathcal{Q}$-filter} of $\mathcal{L}$  if  for all $u,v \in \mathcal{L}$,   $u \vee v \in \mathcal{P}$ and $u \notin \mathcal{Q}$ imply that   $v \in \mathcal{P}$.
\end{definition}
\begin{example} \label{exa1}
\begin{enumerate}
\item Let $\mathcal{L}=\{0,u,v,w,1\}$ with the ralations $u < v, u \wedge w=0$ and $v \vee w=1$. Set $\mathcal={P}=\{1,a\}$ and $\mathcal{Q}=\{1,u,v\}$. Then, $\mathcal{P}$ is a $\mathcal{Q}$-filter of $\mathcal{L}$.
\item Consider the lattice $\mathcal{L} = \{0, u, v, w, 1\}$   in which  $0 \le u \le w \le 1$, $0 \le v \le w \le 1$, $u \vee v = w$, and $u \wedge v = 0$. In the lattice $\mathcal{L}$, $\mathcal{P}=\{1.w\}$, $\mathcal{Q}_1=\{1,w,u\}$ and $\mathcal{Q}_2=\{1,w,v\}$ are filters.   Although  $u \vee v =w\in \mathcal{P}$ and $v \notin \mathcal{Q}_1$, $u \notin  \mathcal{P}$. This implies that $\mathcal{P}$ is not a $\mathcal{Q}_1$-filter of $\mathcal{L}$. Similarly, one  can  check the filter $\mathcal{P}$  is  not a  $\mathcal{Q}_2$-filter  of $\mathcal{L}$.
 \end{enumerate}
\end{example}
\begin{proposition}
If $\mathcal{P}$ is a {\it $\mathcal{Q}$-filter} of $\mathcal{L}$, then $\mathcal{P} \subseteq  \mathcal{Q}$. 
\end{proposition}
\begin{proof}
On the contrary, assume that $\mathcal{P} \nsubseteq  \mathcal{Q}$. Then, there exists $u \in \mathcal{P}$ such that  $u \notin  \mathcal{Q}$. Since $\mathcal{P}$ is a {\it $\mathcal{Q}$-filter} of $\mathcal{L}$, $u =u \vee 0 \in \mathcal{P}$ and $u \notin  \mathcal{Q}$, we get $0 \in \mathcal{P}$, a contradiction. Therefore, we conclude that $\mathcal{P} \subseteq  \mathcal{Q}$. 
\end{proof}
Recall from \cite{Atani2} that an element $u \in \mathcal{L}$ is  the {\it identity join} of $\mathcal{L}$ if there exists $1 \neq v \in \mathcal{L}$ such that $u \vee v =1$. By $\text{Id}(\mathcal{L})$, we mean the set of all identity joins of $\mathcal{L}$. Moreover, 
a lattice $\mathcal{L}$ containing  $1$ is  an {\it $\mathcal{L}$-domain} if for any $u, v \in \mathcal{L}$, $ u \vee v = 1$ implies $u = 1$ or $ v = 1$.
 Consequently, $\mathcal{L}$ is an $\mathcal{L}$-domain if and only if $\{1\}$ is a prime filter of $\mathcal{L}$.
 \begin{proposition} \label{1}
 Let $\mathcal{Q}$ be a filter of $\mathcal{L}$ such that $\mathcal{Q} \neq \{1\}$. The filter $\{1\}$ is a $\mathcal{Q}$-filter of $\mathcal{L}$ if and only if $\text{Id}(\mathcal{L}) \subseteq \mathcal{Q}$. Furthermore,  if $\mathcal{L}$ is an $\mathcal{L}$-domain, then  $\{1\}$ is a $\mathcal{Q}$-filter of $\mathcal{L}$ for every filter $\mathcal{Q}$ of $\mathcal{L}$.
 \end{proposition}
 \begin{proof}
$\Longrightarrow$ Let   $v \in \text{Id}(\mathcal{L})$ but $v \notin  \mathcal{Q}$. Then, there exists $1 \neq u \in \mathcal{L}$ such that $u \vee v =1$. Since  $\{1\}$ is a $\mathcal{Q}$-filter of $\mathcal{L}$ and $v \notin \mathcal{Q}$, we get $u \in \{1\}$, a contradiction. Hence, $\text{Id}(\mathcal{L}) \subseteq \mathcal{Q}$.

$\Longleftarrow$ Let $u \vee v \in \{1\}$ for $u,v \in \mathcal{L}$ and $v \notin \mathcal{Q}$. This implies that $v \notin \text{Id}(\mathcal{L})$. Therefore, we conclude that $u=1$ as $u \vee v =1$. Thus, the filter $\{1\}$ is a $\mathcal{Q}$-filter of $\mathcal{L}$.
 

Now, we show that $``\text{furthermore}"$ statement holds. Let $\mathcal{Q}$ be an arbitrary filter of $\mathcal{L}$. Assume that $u \vee v \in \{1\}$ such that $ v \notin \mathcal{Q}$. Since $\mathcal{L}$ is an $\mathcal{L}$-domain  and $v \neq 1$, we get $u=1$. Consequently, $\{1\}$ is a $\mathcal{Q}$-filter of $\mathcal{L}$.
 \end{proof}
 

\begin{theorem} \label{2}
Let $\mathcal{P}$ and $ \mathcal{Q}$ be   filters  of $\mathcal{L}$ such that  $\mathcal{P} \neq \{1\}$.  The following statements are equivalent: 
\begin{itemize} 
\item[\rm(i)]~ $\mathcal{P}$ is a  $\mathcal{Q}$-filter of $\mathcal{L}$;
\item[\rm(ii)]~  For any two filters $\mathcal{I}$ and $\mathcal{J}$  of $\mathcal{L}$, $\mathcal{I} \vee \mathcal{J} \subseteq \mathcal{P}$ and $\mathcal{I} \nsubseteq \mathcal{Q}$ imply that   $\mathcal{J}  \subseteq \mathcal{P} $.
\end{itemize}
\end{theorem}

\begin{proof}
(i) $\Longrightarrow$ (ii) Let $\mathcal{P}$ be a  $\mathcal{Q}$-filter of $\mathcal{L}$. Assumee that $\mathcal{I} \vee \mathcal{J} \subseteq \mathcal{P}$ for filters  $\mathcal{I}$ and $\mathcal{J}$  of $\mathcal{L}$ such that $\mathcal{I}  \nsubseteq  \mathcal{Q}$. Then, there exists $u \in \mathcal{I}$ such that $ u \notin  \mathcal{Q}$. Now, take any $v \in \mathcal{J}$.  Since  $\mathcal{P}$ is a  $\mathcal{Q}$-filter of $\mathcal{L}$, $ u \vee v \in \mathcal{I} \vee \mathcal{J} \subseteq \mathcal{P}$ and $u \notin \mathcal{Q}$, we get  $v \in \mathcal{P}$ and so $\mathcal{J}  \subseteq \mathcal{P} $.   Hence, (ii) holds.

(ii) $\Longrightarrow$ (i) Let $u \vee v \in \mathcal{P}$ for $u,v \in \mathcal{L}$ and $u \notin \mathcal{Q}$. Put $\mathcal{I}=T(\{u\})$ and $\mathcal{J}=T(\{v\})$.   Since  $\mathcal{I} \vee \mathcal{J} \subseteq \mathcal{P}$ and $\mathcal{I} \nsubseteq \mathcal{Q}$, we obtain $\mathcal{J} \subseteq \mathcal{P}$ or  by the assumption. This implies that $v \in \mathcal{P}$. Thus, $\mathcal{P}$ is a $\mathcal{Q}$-filter of $\mathcal{L}$.
\end{proof}

Let $J(\mathcal{L})$ is the intersection of all maximal filters of $\mathcal{L}$. A proper filter $\mathcal{P}$ of $\mathcal{L}$ is called a {\it $J$-filter} if whenever $u,v \in \mathcal{L}$ with $u \vee v \in \mathcal{P}$ and $v \notin J(\mathcal{L})$, then $u \in \mathcal{P}$ \cite{Atani6}.
\begin{theorem} \label{3}
Let $\mathcal{P} \neq \{1\}$ be   a filter of $\mathcal{L}$.  Then,  $\mathcal{P}$ is a  $J$-filter of $\mathcal{L}$ if and only if  $\mathcal{P}$ is an  $\mathcal{M}$-filter of $\mathcal{L}$ for every maximal filter $\mathcal{M}$ of $\mathcal{L}$.
\end{theorem}
\begin{proof}
$\Longrightarrow$ Let $\mathcal{P}$ be a  $J$-filter of $\mathcal{L}$ and $\mathcal{M}$ be an arbitrary filter of $\mathcal{L}$. Then,   we have $J(\mathcal{L}) \subseteq \mathcal{M}$. Assume that $u \vee v \in \mathcal{P}$ such that $ u \notin \mathcal{M}$. Therefore, we get  $v \in  \mathcal{P}$ as $\mathcal{P}$ is a  $J$-filter of $\mathcal{L}$ and $u \notin  J(\mathcal{L})$. Hence, $\mathcal{P}$ is an  $\mathcal{M}$-filter of $\mathcal{L}$.

$\Longleftarrow$ Assume that $\mathcal{P}$ is an  $\mathcal{M}$-filter of $\mathcal{L}$ for every maximal filter $\mathcal{M}$ of $\mathcal{L}$. Let $u \vee v \in \mathcal{P}$ and $u \notin J(\mathcal{L})$. Then, there is at least a maximal  filter, say $\mathcal{M}_0$,  that do not contain $u$.  By the hypothesis, $\mathcal{P}$ is an  $\mathcal{M}_0$-filter of $\mathcal{L}$. This implies that   $v \in \mathcal{P}$ as $u \vee v \in \mathcal{P}$ and $u \notin \mathcal{M}_0$. Thus, $\mathcal{P}$ is a  $J$-filter of $\mathcal{L}$.
\end{proof}

\begin{thebibliography}{99}
\bibitem[Atani2]{Atani2} S.E. Atani, {\it $S$-prime property in lattices}, Mathematica, {\bf 66} (89) (2024) 168–187.
\bibitem[Atani6]{Atani6} S.E. Atani, {\it $J$-Filters of bounded lattices}, Eur. J. Math. Appl., {\bf 4}(12) (2024) 1-13.
\end{thebibliography}
\end{document}
